% sec-03: who is responsible for what.  Figure 7 (graphviz-type clusters),
% Table 15.
\section{Who Is Responsible for What}

A site agreement is blocked more often by an unanswered responsibility question
than by a scientific one. This section states the split as it would stand, marks
everything conditional as conditional, and lists the four questions that are open.

\begin{appfloat}[!tb]
\begin{appfig}[graphviz-type / three dashed clusters with edges only between them]
% Three clusters on a 4.75cm pitch; nodes on a 0.92cm pitch.  Every edge runs
% between clusters and none runs inside one, which is what keeps the clusters
% readable.  Developer edges are dashed: they are conditional on an agreement
% that does not exist.
% Node names are plain letters and digits.  A generated name carrying a minus or
% a dot would be read as a subtraction or as a node-and-anchor pair.
\node[gvboxs,text width=26mm] (sa) at (0,0)     {Dose assignment};
\node[gvboxs,text width=26mm] (sb) at (0,-0.92) {Dose-limiting toxicity adjudication};
\node[gvboxs,text width=26mm] (sc) at (0,-1.84) {Analysis execution and unblinding};
\node[gvboxs,text width=26mm] (sd) at (0,-2.76) {Investigator financial disclosure};
\node[gvboxm,text width=26mm] (pa) at (4.75,0)     {Protocol authorship and amendment};
\node[gvboxm,text width=26mm] (pb) at (4.75,-0.92) {Analysis plan, pre-specified};
\node[gvboxm,text width=26mm] (pc) at (4.75,-1.84) {Safety reporting and annual reports};
\node[gvboxm,text width=26mm] (pd) at (4.75,-2.76) {Insurance and subject injury};
\node[gvboxg2,text width=26mm] (da) at (9.50,0)     {Drug supply};
\node[gvboxg2,text width=26mm] (db) at (9.50,-0.92) {Letter of authorization};
\node[gvboxg2,text width=26mm] (dc) at (9.50,-1.84) {Pharmacovigilance requirements};
\node[gvctitle,anchor=south west] (ct1) at (-1.28,0.62) {Site, under an agreement that does not yet exist};
\node[gvctitle,anchor=south west] (ct2) at (3.47,0.62) {Sponsor, ChemicalQDevice};
\node[gvctitle,anchor=south west] (ct3) at (8.22,0.62) {Developer, no agreement exists};
\node[gvcluster,fit={(sa)(sd)(ct1)},inner sep=6pt] {};
\node[gvcluster,fit={(pa)(pd)(ct2)},inner sep=6pt] {};
\node[gvcluster2,fit={(da)(dc)(ct3)},inner sep=6pt] {};
\draw[gvedgeb] (pa.west) -- node[above,font=\tiny] {site concurs} (sa.east);
\draw[gvedgeb] (pb.west) -- node[above,font=\tiny] {pre-specified} (sc.east);
\draw[gvedged] (da.west) -- node[above,font=\tiny] {if agreed} (pa.east);
\draw[gvedged] (db.west) -- node[above,font=\tiny] {if agreed} (pc.east);
\draw[gvedgeg] (sd.south) -- ++(0,-0.50)
  -- node[midway,below,font=\tiny] {collected and maintained} (4.75,-3.26) -- (pd.south);
\node[font=\tiny\itshape,text=fundgray,anchor=north west,text width=118mm]
  at (-1.28,-4.05) {Dashed edges are conditional on an agreement that does not
  exist. No approach beyond a first feasibility letter has been made to the
  developer, and no institution has agreed to anything.};
\end{appfig}
\vspace{-0.60cm}
\figcaption{Figure 7. Site, sponsor and developer obligations as three disjoint sets, with\\
edges only between them and dashed lines wherever no agreement yet exists.}
\end{appfloat}

\subsection{Why the site holds the investigator role}

The chief executive is chief executive of the sponsor and holds the whole of it.
He is not a clinical investigator on this trial, will not be one, and is nowhere
described as one. If he held both roles, all four of the 21 CFR 54.2 disclosable
interests would fire at once: compensation influenced by the study outcome,
significant equity in a non-publicly traded sponsor, a proprietary interest in
the workflow assets, and significant payments of other sorts
\cite{ecfr2026part54,fdafindisclosure2013}.

Because all four would fire, the separation cannot be partial. A Phase 1 3+3
escalation carries no randomization, so the boundary binds on dose assignment,
dose-limiting toxicity adjudication, and the analysis plan instead, and the
wording is carried forward so that it still binds when a Phase 2 successor
introduces randomization \cite{phase2proto}.

\subsection{The three corrections every letter carries}

\begin{apptable}
\begin{tabularx}{\textwidth}{>{\raggedright\arraybackslash}p{0.5cm} >{\raggedright\arraybackslash}p{5.2cm} >{\raggedright\arraybackslash}X}
\headrow \thead{\#} & \thead{The correction} & \thead{Why it travels with every communication}\\
\midrule
1 & Daraxonrasib is not described as first in human & It is approved in the metastatic setting and was in Phase 3 evaluation before that. The supportable claim concerns the integrated surgical and advisory workflow \\
2 & The robotic configuration is specified at the site agreement and not before & No site can assess feasibility or regulatory status without an identifiable platform, and asserting one this company does not have would be worse than asserting none \\
3 & No institution is a partner, sponsor, site, or endorser & No agreement of any kind exists, and a single sentence implying otherwise would be found by the institution named in it \\
\end{tabularx}
\end{apptable}
\vspace{-0.60cm}
\tabcap{Table 15. The three positioning corrections carried into every letter,\\
with the reason each one travels rather than being stated once and filed.}

\subsection{The four questions this section leaves open}

Who holds the investigational new drug application, the sponsor or the developer.
Whether a site physician could serve as sponsor-investigator, and whether that
would be better or worse than the current split. Who performs safety reporting,
monitoring, pharmacovigilance, drug accountability and annual reports. And who
pays for insurance, indemnification, subject injury, and device and software
validation \cite{ecfr2026part312records}.

All four are asked in the fifth letter of this day. None is answered here,
because answering them in a document the institution has not seen would be
answering them alone.

\subsection{The confidentiality route, asked rather than assumed}

No confidential engineering detail travels before a confidentiality agreement is
in place, and the route that agreement should follow is a question in the fifth
letter of this day rather than an assumption in the first. The project is
investigator-authored, seeks federal support, has a for-profit applicant, and may
require drug support from a third party, and each of those facts can change which
institutional office should hold the agreement.

Asking is cheaper than guessing. An agreement executed through the wrong office
is renegotiated through the right one, and the renegotiation happens at the point
where the detail is actually needed, which is the worst moment to discover it.

\subsection{What a reader should take from this section}

That the responsibility split is written down before it is negotiated, that
everything conditional is marked conditional in the figure as well as in the
text, and that the four open questions are named rather than papered over. A
site that receives a document asserting a split it never agreed to reads a
company that will be difficult later. A site that receives one marking its own
gaps reads a company that has thought about the agreement it is asking for.

